% Inline Autogyro Mechanical Assembly — Cross-Section Comparison
% Compile: pdflatex inline_autogyro_mech.tex
\documentclass[tikz,border=14pt]{standalone}
\usepackage{amsmath}
\usepackage{amssymb}
\usepackage{tikz}
\usetikzlibrary{arrows.meta,calc,patterns,decorations.pathreplacing,positioning,shapes.geometric}

\begin{document}
\begin{tikzpicture}[
    scale=1.0,
    % Line styles
    dyneema/.style={draw=blue!50!black, line width=2.0pt},
    tube/.style={draw=gray!60, fill=gray!15, line width=1.5pt},
    metal/.style={draw=gray!40, fill=gray!25, line width=1.0pt},
    bearing/.style={draw=gray!50, fill=gray!20, line width=1.0pt},
    rotor/.style={draw=red!60!black, fill=red!8, line width=1.5pt},
    blade/.style={draw=red!50!black, fill=red!5, line width=1.2pt},
    force/.style={-{Stealth[scale=1.8]}, line width=2.0pt},
    annot/.style={font=\footnotesize\sffamily, fill=white, fill opacity=0.9,
                  text opacity=1, inner sep=1.5pt, rounded corners=1pt},
    dimension/.style={latex-latex, thin, draw=gray!60},
    label/.style={font=\small\sffamily\bfseries},
]

% =============================================================
% APPROACH A: ANGLED BEARING SEAT (left)
% =============================================================

\def\tilt{12}    % disk tilt angle (degrees)
\def\lineangle{55} % line elevation from horizontal

\begin{scope}[local bounding box=approachA, shift={(0,0)}]

% Draw line direction
\coordinate (Ltop) at (0, 7.5);
\coordinate (Lbot) at (0, 0);
\coordinate (Ldir) at (0, 1);

% === Dyneema line ===
\draw[dyneema] (Lbot) -- (Ltop);
\node[annot] at (-0.4, 6.8) {$T_\text{above}$};
\node[annot] at (-0.4, 0.5) {$T_\text{below}$};

% === Compression tube ===
% Tube extends from below the bearing down to the empennage attachment
\fill[tube] (-0.25, 1.5) rectangle (0.25, 5.0);
\draw[tube] (-0.25, 1.5) -- (-0.25, 5.0) -- (0.25, 5.0) -- (0.25, 1.5);
% Inner hole for line
\fill[white] (-0.1, 1.5) rectangle (0.1, 5.0);
\draw[thin] (-0.1, 1.5) -- (-0.1, 5.0);
\draw[thin] (0.1, 1.5) -- (0.1, 5.0);

\node[annot, right=0.3cm] at (0.25, 3.2) {Tube\\(compression)};

% === Angled bearing seat at tube top ===
% Tube top is cut at tilt angle
\begin{scope}
    \clip (-1.5, 4.5) rectangle (1.5, 5.5);
    \fill[tube, rotate around={\tilt:(0,5.0)}] (-0.35, 4.95) rectangle (0.35, 5.15);
\end{scope}

% Bearing housing
\fill[bearing, rotate around={\tilt:(0,5.0)}]
    (-0.4, 5.0) rectangle (0.4, 5.3);
\draw[bearing, rotate around={\tilt:(0,5.0)}]
    (-0.4, 5.0) -- (-0.4, 5.3) -- (0.4, 5.3) -- (0.4, 5.0);

% Bearing race markers (balls)
\foreach \x in {-0.3, -0.15, 0, 0.15, 0.3} {
    \fill[gray!60, rotate around={\tilt:(0,5.0)}]
        (\x, 5.15) circle (0.04);
}

% Bearing seat angle annotation
\draw[dimension] (0.6, 5.0) -- ++({0.6*sin(\tilt)}, {-0.6*cos(\tilt)});
\node[annot, right=0.1cm] at (0.9, 4.7) {$\delta$};

\node[annot] at (-1.3, 5.3) {Thrust\\bearing};

% === Rotor hub (hollow, around line) ===
% The hub sits on the thrust bearing and autorotates
\fill[metal, rotate around={\tilt:(0,5.0)}]
    (-0.35, 5.3) rectangle (0.35, 5.8);
\fill[white, rotate around={\tilt:(0,5.0)}]
    (-0.08, 5.25) rectangle (0.08, 5.85);

\node[annot, right=0.3cm] at (0.6, 6.0) {Hollow\\hub};

% === Rotor disk (annulus around hub) ===
\fill[rotor, rotate around={\tilt:(0,5.0)}]
    (-1.5, 5.4) rectangle (1.5, 5.6);
\fill[white, rotate around={\tilt:(0,5.0)}]
    (-0.35, 5.35) rectangle (0.35, 5.65);

\node[label] at (-2.0, 5.8) {Rotor disk};

% === Blades (shown in cross-section) ===
\fill[blade, rotate around={\tilt:(0,5.5)}]
    (-2.3, 5.45) rectangle (-1.5, 5.55);
\fill[blade, rotate around={\tilt:(0,5.5)}]
    (1.5, 5.45) rectangle (2.3, 5.55);

\node[annot] at (-3.2, 5.8) {Blade};

% === Swashplate (below rotor, around tube) ===
% Stationary (non-rotating) ring attached to tube
\fill[metal] (-0.45, 4.2) rectangle (0.45, 4.5);
% Rotating ring attached to rotor hub
\fill[bearing] (-0.4, 4.5) rectangle (0.4, 4.7);
% Links connecting to tube
\draw[force, gray] (-0.35, 4.35) -- (-0.35, 3.8);
\draw[force, gray] (0.35, 4.35) -- (0.35, 3.8);

\node[annot, right=0.4cm] at (0.6, 4.3) {Swashplate\\(non-rotating)};

% === Actuators (servos, mounted on tube) ===
\fill[metal] (-0.2, 3.3) rectangle (0.1, 3.8);
\fill[metal] (0.1, 3.3) rectangle (0.4, 3.8);
\node[annot, right=0.3cm] at (0.5, 3.5) {Actuators};

% === Empennage (tail boom + stabilizer) ===
% Mounted on bottom of tube, extends downwind
\fill[metal] (0.25, 1.8) rectangle (2.2, 2.2);
\node[annot] at (1.8, 2.5) {Tail boom};

% Horizontal stabilizer
\fill[metal] (1.5, 2.0) rectangle (2.2, 3.0);
\node[annot] at (2.8, 2.5) {H-stab};

% Vertical fin
\fill[metal] (1.5, 2.5) rectangle (2.0, 3.5);
\node[annot] at (2.8, 3.0) {V-fin};

% === Forces ===
% Lift force vector from rotor
\draw[force, green!50!black, very thick]
    (0, 5.5) -- ++({1.2*sin(\tilt)}, {1.2*cos(\tilt)})
    node[right, annotplain, green!50!black] {$F_\text{lift}$};

% Line tension above (downward through bearing)
\draw[force, purple!50!black, very thick]
    (0, 5.8) -- (0, 5.0)
    node[left, annotplain, purple!50!black] {$T_\text{above}$};

% Compression in tube
\draw[force, purple!50!black, line width=3pt]
    (0, 5.0) -- (0, 1.5);

% === Title ===
\node[label, font=\large] at (0, -1.0) {\textbf{Approach A: Angled bearing seat}};

% Features box
\node[draw=black!30, fill=white, rounded corners=3pt, inner sep=6pt,
      anchor=north west, font=\scriptsize\sffamily, align=left]
    at (-2.5, 7.8) {
    \textbf{Key features:}\\
    $\bullet$ Tube top cut at tilt angle\\
    $\bullet$ Thrust bearing on angled face\\
    $\bullet$ Pure thrust on bearing\\
    $\bullet$ Hub autorotates around line\\
    $\bullet$ Swashplate below rotor\\
    $\bullet$ Empennage trims tilt
};

\end{scope}

% =============================================================
% APPROACH B: FLAT BEARING + TILTED HUB (middle)
% =============================================================

\begin{scope}[local bounding box=approachB, shift={(7,0)}]

\coordinate (Ltop) at (0, 7.5);
\coordinate (Lbot) at (0, 0);

% Dyneema line
\draw[dyneema] (Lbot) -- (Ltop);

% Tube (shorter, flat top)
\fill[tube] (-0.25, 1.5) rectangle (0.25, 5.0);
\fill[white] (-0.1, 1.5) rectangle (0.1, 5.0);
\draw[tube] (-0.25, 1.5) -- (-0.25, 5.0) -- (0.25, 5.0) -- (0.25, 1.5);

\node[annot, right=0.3cm] at (0.25, 3.2) {Tube\\(compression)};

% Flat thrust bearing
\fill[bearing] (-0.45, 5.0) rectangle (0.45, 5.25);
\foreach \x in {-0.3, -0.15, 0, 0.15, 0.3} {
    \fill[gray!60] (\x, 5.12) circle (0.04);
}

% Tilted hub — wedge on bottom
\fill[metal] (-0.35, 5.25) -- (0.35, 5.25) -- (0.35, 5.6) -- (-0.35, 5.35) -- cycle;
% Hub body above
\fill[metal] (-0.35, 5.35) rectangle (0.35, 5.85);
\fill[white] (-0.08, 5.2) rectangle (0.08, 5.9);

% Rotor disk (tilted)
\fill[rotor, rotate around={\tilt:(0,5.5)}]
    (-1.5, 5.4) rectangle (1.5, 5.6);

% Blades
\fill[blade, rotate around={\tilt:(0,5.5)}]
    (-2.3, 5.45) rectangle (-1.5, 5.55);
\fill[blade, rotate around={\tilt:(0,5.5)}]
    (1.5, 5.45) rectangle (2.3, 5.55);

% Off-axis moment annotation
\draw[force, red!60!black, very thick, ->]
    (0.6, 5.25) arc[start angle=0, end angle=30, radius=0.6cm];
\node[annot, red!60!black] at (1.4, 5.5) {Moment\\load};

% Empennage
\fill[metal] (0.25, 1.8) rectangle (2.2, 2.2);
\fill[metal] (1.5, 2.0) rectangle (2.2, 3.0);
\fill[metal] (1.5, 2.5) rectangle (2.0, 3.5);

% Title
\node[label, font=\large] at (0, -1.0) {\textbf{Approach B: Flat bearing + tilted hub}};

% Issues box
\node[draw=red!40, fill=red!5, rounded corners=3pt, inner sep=6pt,
      anchor=north west, font=\scriptsize\sffamily, align=left]
    at (-2.5, 7.8) {
    \textbf{Issues:}\\
    $\bullet$ Bearing sees off-axis moment\\
    $\bullet$ Uneven wear on bearing\\
    $\bullet$ Tilted hub wedge is a\\\quad stress concentrator\\
    $\bullet$ Simpler to machine though
};

\end{scope}

% =============================================================
% APPROACH C: U-JOINT + FLOATING TILT (right)
% =============================================================

\begin{scope}[local bounding box=approachC, shift={(14,0)}]

\coordinate (Ltop) at (0, 7.5);
\coordinate (Lbot) at (0, 0);

% Dyneema line
\draw[dyneema] (Lbot) -- (Ltop);

% Tube
\fill[tube] (-0.25, 1.5) rectangle (0.25, 4.5);
\fill[white] (-0.1, 1.5) rectangle (0.1, 4.5);
\draw[tube] (-0.25, 1.5) -- (-0.25, 4.5) -- (0.25, 4.5) -- (0.25, 1.5);

% U-joint yoke (lower, on tube)
\fill[metal] (-0.4, 4.5) rectangle (-0.2, 4.9);
\fill[metal] (0.2, 4.5) rectangle (0.4, 4.9);
\draw[metal] (-0.4, 4.7) -- (-0.2, 4.7);
\draw[metal] (0.2, 4.7) -- (0.4, 4.7);

% U-joint cross pin
\fill[metal] (-0.35, 4.7) circle (0.06);

% U-joint yoke (upper, on bearing housing)
\fill[metal, rotate around={\tilt:(0,4.9)}]
    (-0.4, 4.8) rectangle (-0.2, 5.2);
\fill[metal, rotate around={\tilt:(0,4.9)}]
    (0.2, 4.8) rectangle (0.4, 5.2);

% Thrust bearing on top of U-joint
\fill[bearing, rotate around={\tilt:(0,5.0)}]
    (-0.4, 5.0) rectangle (0.4, 5.25);
\foreach \x in {-0.3, -0.15, 0, 0.15, 0.3} {
    \fill[gray!60, rotate around={\tilt:(0,5.0)}]
        (\x, 5.12) circle (0.04);
}

% Hub
\fill[metal, rotate around={\tilt:(0,5.0)}]
    (-0.35, 5.25) rectangle (0.35, 5.8);
\fill[white, rotate around={\tilt:(0,5.0)}]
    (-0.08, 5.2) rectangle (0.08, 5.85);

% Rotor disk
\fill[rotor, rotate around={\tilt:(0,5.5)}]
    (-1.5, 5.4) rectangle (1.5, 5.6);

% Blades
\fill[blade, rotate around={\tilt:(0,5.5)}]
    (-2.3, 5.45) rectangle (-1.5, 5.55);
\fill[blade, rotate around={\tilt:(0,5.5)}]
    (1.5, 5.45) rectangle (2.3, 5.55);

% Empennage
\fill[metal] (0.25, 1.8) rectangle (2.2, 2.2);
\fill[metal] (1.5, 2.0) rectangle (2.2, 3.0);
\fill[metal] (1.5, 2.5) rectangle (2.0, 3.5);

% U-joint annotation
\node[annot] at (-1.2, 4.7) {U-joint};

% Title
\node[label, font=\large] at (0, -1.0) {\textbf{Approach C: U-joint + floating tilt}};

% Issues box
\node[draw=orange!60, fill=orange!5, rounded corners=3pt, inner sep=6pt,
      anchor=north west, font=\scriptsize\sffamily, align=left]
    at (-2.5, 7.8) {
    \textbf{Issues:}\\
    $\bullet$ 2 rotational DOF — pitch + roll\\
    $\bullet$ Empennage must stabilize both\\
    $\bullet$ Bearing sees pure thrust (good)\\
    $\bullet$ U-joint adds weight + complexity\\
    $\bullet$ Wind gusts could excite oscillation
};

\end{scope}

% =============================================================
% TITLE
% =============================================================

\node[label, font=\Large] at (7, 8.5)
    {Inline Autogyro Kite — Mechanical Assembly Approaches};

\end{tikzpicture}
\end{document}
